% Generated by `taskgen course examples` from number_theory.solovay_strassen (seed 1).
% Do not edit: change examples.yaml and regenerate.
\begin{example}[тест Соловея--Штрассена]\label{ex:solovay-strassen-prime}
С помощью теста Соловея-Штрассена определить, является ли простым число $N = 457$. Требуется дать ответ с вероятностью ошибки не~более~$\frac14$.

\par\noindent\textit{Решение.} Для простого $N$ при любом $a$, $\gcd(a, N) = 1$: $a^{(N-1)/2} \equiv \left(\frac{a}{N}\right)\ (\textrm{mod}\ N)$; для составного $N$ это выполнено не более чем для половины $a$, поэтому нужно два раунда ($(1/2)^2 = 1/4$). $(N-1)/2 = 228$.\par\noindent $a = 2$: $\left(\frac{2}{457}\right) = \left(\frac{1}{457}\right) = 1$; $2^{228} \equiv 1$ --- совпадает (раунд пройден).\par\noindent $a = 3$: $\left(\frac{3}{457}\right) = \left(\frac{457}{3}\right) = \left(\frac{1}{3}\right) = 1$; $3^{228} \equiv 1$ --- совпадает (раунд пройден).\par\noindent Пройдено 2 раунда, вероятность ошибки не более $\frac14$: $N = 457$ --- простое.
\end{example}
